
\documentclass[aspectratio=169]{beamer}

\mode<presentation>
\usetheme{Boadilla}
\definecolor{NTHU1}{RGB}{127,16,132}
\definecolor{NTHU2}{RGB}{228,210,231}
\definecolor{blue}{RGB}{30,90,205}
\definecolor{red}{RGB}{213,94,0}
\definecolor{green}{RGB}{0,128,0}
\definecolor{charcoalgray}{RGB}{108,104,104}
\setbeamercolor{title}{fg=NTHU1}
\setbeamercolor{frametitle}{fg=NTHU1}
\setbeamercolor{block title}{bg=NTHU1, fg=white}
\setbeamercolor{block body}{bg=NTHU2}
\setbeamercolor{structure}{fg=NTHU1}
\setbeamercolor{item projected}{fg=white}
\setbeamercolor{item}{fg=NTHU1}
\setbeamercolor{subitem}{fg=NTHU1}
\setbeamercolor{section in toc}{fg=NTHU1}
\setbeamercolor{description item}{fg=NTHU1}
\setbeamercolor{caption name}{fg=NTHU1}
\setbeamercolor{button}{bg=NTHU1, fg=white}
\setbeamercolor{caption name}{fg=NTHU1}
\usepackage{graphics}
\usepackage{tikz}
\usepackage{amsmath}
\usepackage{bbm}
\usetikzlibrary{decorations.pathreplacing}
\usepackage{geometry}
\usepackage{booktabs}
\usepackage{bookmark}
\usepackage{multirow, makecell}
\usepackage{float}
\usepackage{fancyvrb}
\usepackage{caption}
\captionsetup[figure]{singlelinecheck = false, format= hang, justification=centering}
\captionsetup[subfigure]{singlelinecheck = false, format= hang, justification=raggedright}
\usepackage{subcaption}
\usepackage{adjustbox}
\usepackage{hyperref}
\usepackage{threeparttable}
\usepackage{appendixnumberbeamer} 
\usepackage[T1]{fontenc}
\usepackage[default]{lato} 
\usepackage{smartdiagram}
\smartdiagramset{
  back arrow disabled = true,
  module minimum width=1cm,
  module x sep = 3,
  uniform arrow color=true,
}
\usepackage{xcolor}
\usepackage{colortbl}
	\definecolor{light-gray}{gray}{0.99}

\newenvironment{wideitemize}{\itemize\addtolength{\itemsep}{10pt}}{\enditemize}
\newenvironment{wideenumerate}{\enumerate\addtolength{\itemsep}{10pt}}{\endenumerate}
\newenvironment{widedescription}{\description\addtolength{\itemsep}{10pt}}{\enddescription}
\hypersetup{
colorlinks=true,
linkcolor=NTHU1,
filecolor=green, 
urlcolor=blue,
}
\beamertemplatenavigationsymbolsempty
\setbeamercolor{author in head/foot}{bg=white, fg=NTHU1}
\setbeamercolor{title in head/foot}{bg=white, fg=NTHU1}
\setbeamercolor{date in head/foot}{bg=white, fg=NTHU1}
\setbeamercolor{section in head/foot}{bg=white, fg=NTHU1}
\setbeamercolor{headline}{bg=white}
\setbeamertemplate{footline}{
    \leavevmode%
    \hbox{%
        \begin{beamercolorbox}[wd=.333333\paperwidth,ht=2.25ex,dp=1ex,center]{date in head/foot}%
            \usebeamerfont{date in head/foot}%\insertdate
        \end{beamercolorbox}%
        \begin{beamercolorbox}[wd=.333333\paperwidth,ht=2.25ex,dp=1ex,center]{title in head/foot}%
            \usebeamerfont{title in head/foot}\insertshorttitle
        \end{beamercolorbox}%
        \begin{beamercolorbox}[wd=.333333\paperwidth,ht=2.25ex,dp=1ex,right]{date in head/foot}%
            \usebeamerfont{date in head/foot}\insertframenumber{}\hspace*{2em}
        \end{beamercolorbox}}%
        \vskip0pt%
    }
    %% May need to script the introduction!!!!!!
%\setbeamercolor{page number in head/foot}{fg=NTHU1}
\setbeamertemplate{section in toc}[sections numbered]
\setbeamertemplate{subsection in toc}{\leavevmode\leftskip=3em\rlap{\hskip-1.75em\inserttocsectionnumber.\inserttocsubsectionnumber}
\inserttocsubsection\par}
\setbeamerfont{subsection in toc}{size=\footnotesize}
%\setbeamertemplate{headline}{%
  %\begin{beamercolorbox}[ht=5.5ex]{section in head/foot}
    %\vskip2pt\insertnavigation{0.33\paperwidth}\vskip2pt
  %\end{beamercolorbox}%
%}
\newenvironment{transitionframe}{\setbeamercolor{background canvas}{bg=white}\setbeamertemplate{footline}{} \begin{frame}[noframenumbering]}{\end{frame}}


\makeatletter
\let\@@magyar@captionfix\relax
\makeatother


\title[]{Public Economics and Development} % Change this regularly
\subtitle[]{Week 8: Evolution of tax structures in developing countries}
\author[Lee]{Seung-hun Lee }
\institute[]{Taipei School of Economics\\ National Tsing Hua University}

\date[]{April 17th, 2026}

\begin{document}
\begin{transitionframe}
\titlepage
\end{transitionframe}


%%% Color slides for section headers: Use for colloquium version (The ones bewteen \iffals and \fi)

\section{Overview setup of the topic}

\begin{frame}
\frametitle{Challenges in tax administration: Information} 
Information on taxable transactions are cricial for imposing tax
\begin{itemize}\setlength\itemsep{3pt}
\item Having a firm registered
\item Income sources with traceable sources
\item Formal transactions with verifiable paper trails
\end{itemize}\medskip
Modern economies implement technologies and systems gathering these information
\begin{itemize}\setlength\itemsep{3pt}
\item Statistical offices \& population registers compile info on various economic activities
\item Employee-employer relationships enable accurate tracking of income
\begin{itemize}
\item Third-party information (not self-declared information) used as source
\end{itemize}
\item Value-added tax (VAT) creates a paper trail between different levels of supply chain
 \begin{itemize}\setlength\itemsep{3pt}
\item VAT charged on sales - input for each firms
\item Downstream firms have incentive to ask their suppliers for a receipt
\item Since the receipt provides evidence for input costs, which can lower their tax burden
\end{itemize}
\end{itemize}
\end{frame}




\begin{frame}
\frametitle{Rise of informational capacity and taxation over time} 
\begin{figure}
  \includegraphics[keepaspectratio, width=0.7\textwidth]{fig1.png}
\end{figure}
\end{frame}

\begin{frame}
\frametitle{Lack of tax administration capacity in developing countries} 
Limited access to necessary information
\begin{itemize}\setlength\itemsep{3pt}
\item Large share of economic activities are taking place in informal markets (last week)
\item Usage of third party information is limited
\begin{itemize}\setlength\itemsep{3pt}
  \item Higher share of self-employed individuals (who self-declare income)
\item Less transactions based on digital tools such as credit cards (cash do not leave records)
\end{itemize}
\item Tax agents are often under-trained or under-monitored
\end{itemize}\medskip
What differences does it lead to
\begin{itemize}\setlength\itemsep{3pt}
  \item Generally lower share of taxes relative to economy
  \item High share of trade taxes such as tariffs (easier to monitor)
\item Collect more taxes from firms than individuals (fewer firms than individuals)
\end{itemize}
\end{frame}

\begin{frame}
\frametitle{Taxes in OECD vs non-OECD countries} 
\begin{figure}
  \includegraphics[keepaspectratio, width=1\textwidth]{fig2.png}
\end{figure}
\end{frame}


\begin{frame}
\frametitle{New developments in tax technology} 
What is being tried in developing countries
\begin{itemize}\setlength\itemsep{3pt}
\item Expansion of digital money and banking system for transparent transactions
\item Increases in digitalized identifiaction of properties and individuals
\item E-filing: Making compliance to taxation easier
\end{itemize}\medskip
Some new studies:
\begin{itemize}\setlength\itemsep{3pt}
  \item Effects of the assignment of tax collectors
  \item Effects of incorporating local knowledge in taxation
  \item Implmentation of various tax campaigns to increase compliance from taxpayers
\end{itemize}
\end{frame}

\begin{transitionframe}
\frametitle{Roadmap for today}
\tableofcontents
\end{transitionframe}








\section{Evolution of economic structure and taxation}
\subsection{Oyebola, Tourek (2024): How Can Lower-Income Countries Collect More Taxes? The Role of Technology, Tax Agents and Politics}
\begin{transitionframe}
\begin{center}
  \begin{huge}
    \textcolor{NTHU1}{%
Oyebola, Tourek (2024):\\ How Can Lower-Income Countries Collect More Taxes? The Role of Technology, Tax Agents and Politics
    }
  \end{huge}
\end{center}
\end{transitionframe}


\begin{frame}
\frametitle{Q: How do you improve tax system in developing countries? } 
Increasing tax revenues is a major goal in many low/middle income countries
\begin{itemize}\setlength\itemsep{3pt}
\item Generally have lower share of their GDP as tax revenues
\item Low taxation is a symptom of economic and institutional constraints
\begin{itemize}
\item Tax base not large enough, lack of constraint on redistribution/taxes
\end{itemize}
\item Does economic growth automatically expand tax capacity? Not exactly
\begin{itemize}
\item Separate investments in improving the tax system required
\end{itemize}
\end{itemize}\medskip
This paper reviews recent developments in tax administration in developing countries
\begin{itemize}\setlength\itemsep{3pt}
\item Identifying taxpayers (expanding tax registries and databases)
\item Detecting tax burden (usage of third-party information)
\item Collect taxes to state coffers (Billing systems, legal measures for noncompliance)
 \end{itemize}
\end{frame}

\begin{frame}
\frametitle{Technology in tax collection} 
Identification: Cost of linking identity info across government agencies and registries$\Downarrow$

\begin{itemize}\setlength\itemsep{3pt}
\item National ID (Ghana Card), Digital mapping of properties (Liberia)
\item ID information may need to be combined with other measures (Okunogbe 2023)
\end{itemize}\medskip
Detection: Collect third-party info (declaration of tax liabilities by third person)
\begin{itemize}\setlength\itemsep{3pt}
\item `Digital trails' from value-added taxes and employment taxes
\item Requires countries to cross-check information using multiple sources
\item Need to close international loopholes
\end{itemize}\medskip
Collection: Increases in electronic filing and mobile payment options
\begin{itemize}\setlength\itemsep{3pt}
\item Reduces cost of monitoring and Compliance
\item May disadvantage those less familiar with electronics \\ (rural, less-educated, less-established firms)
\end{itemize}\medskip

\end{frame}

\begin{frame}
\frametitle{Cross-country evidence on technology and collection} 
\begin{figure}
  \includegraphics[keepaspectratio, width=0.6\textwidth]{ot2024_f1.png}
\end{figure}
\end{frame}

\begin{frame}
\frametitle{Role of tax collectors} 
What determines the effectiveness of officials
\begin{itemize}\setlength\itemsep{3pt}
\item Deployment: low citizen-to-staff ratio and appropriate allocation
\begin{itemize}\setlength\itemsep{3pt}
\item Low burden per staff
\item Specialization: Matching staff depending on tax types and taxpayer traits 
\end{itemize}
\item Incentives: Tie staff revenues to tax revenue generations
 \begin{itemize}\setlength\itemsep{3pt}
\item Risk: Could increase scope for corruption
\item Monitoring staff with audits crucial 
\end{itemize}
\end{itemize}\medskip
Interaction with technology
\begin{itemize}\setlength\itemsep{3pt}
\item Substitute? Obviate need for in-person Interaction
\item Complement? Help expedite tedious processes and collecting contextual knowledge
\item Balancing the two forces is key (and unresolved question)
\item Resistance to adopting the technological change
\end{itemize}\medskip
\end{frame}


\begin{frame}
\frametitle{Political incentives for taxation} 
How do political incentives determine tax collection
\begin{itemize}\setlength\itemsep{3pt}
\item Tax$\Uparrow$ $\to$ fund public services, social programs, infrastructure $\to$ support re-election
\item Tax$\Uparrow$ $\to$ Citizen demand, expectation for accountability $\Uparrow$ $\to$ not help that much
\end{itemize} \medskip
What are the factors shaping success of tax collection?
\begin{itemize}
\item Alternative revenue: Dampen motivation for Taxation
\item Political competition: Competition for better governance (through public services)
\item Availability of tech: Reduce cost of collection and resistance from taxpayers
\end{itemize}
\end{frame}

\begin{frame}
\frametitle{Expanding our knowledge} 
Findings from other relevant studies
\begin{itemize}\setlength\itemsep{3pt}
\item Tax revenue and compliance increases with increased tech and enforcement
\item Many results are from small-scale local experiments: Generalizable?
\end{itemize}\medskip
What remains?
\begin{itemize}\setlength\itemsep{3pt}
\item How to make technology and personnel complementary?
\item How to tax new commodities (digital activities, digital currencies)?
\item How to increase political support for taxation?
\end{itemize}
\end{frame}


\subsection{Jensen (2022): Employment Structure and the Rise of the Modern Tax System}
\begin{transitionframe}
\begin{center}
  \begin{huge}
    \textcolor{NTHU1}{%
Jensen (2022):\\ Employment Structure and the Rise of the Modern Tax System
    }
  \end{huge}
\end{center}
\end{transitionframe}


\begin{frame}
\frametitle{Q: How do modern tax systems rise?}
Modern tax systems are characterized by increased direct taxes (income tax)
\begin{itemize}\setlength\itemsep{3pt}
\item Requires high collection capacities 
\item Changes in the structure of employment may explain increase in collection
\begin{itemize}\setlength\itemsep{3pt}
\item Enforcement in employer-employee structure $>$ Self-employed
\item Existence of third-party information (employer, in case of income tax)
\end{itemize}
\end{itemize}\medskip
This paper uncovers new facts on tax capacity with microdata across countries and time
\begin{itemize}\setlength\itemsep{3pt}
\item Employee share increases over income distribution \& further down the income distribution as country develops 
\item Exemption threshold gradually decreases in income distribution over development
\item Threshold is located such that employee share on tax base remains constant and high
\end{itemize}
\end{frame}

\begin{frame}
\frametitle{Data, definition, and calculation strategy} 
Data source and definitions
\begin{itemize}\setlength\itemsep{3pt}
\item Household surveys (100 countries)
\begin{itemize}\setlength\itemsep{3pt}
\item Employment type of economically active population (EAP)
\item Information on taxable earnings (employment / self-employment / capital / other income)
\item For US (1870-2010) and Mexico (1960-2010), compile long run time series 
\end{itemize}
\end{itemize}\medskip
Definition of key terms and strategies for calculation
\begin{itemize}\setlength\itemsep{3pt}
\item Employee: if work type generates third-party information usable in tax enforcement \\ (Otherwise, self-employment)
\begin{itemize}\setlength\itemsep{3pt}
\item Self-employed: casual laborers, own account workers, workers w/o contracts
\item Calculate share of employees in EAP across country's income distribution (in deciles)
\end{itemize}
\item Threshold: Minimum level of gross income above which individuals have to pay taxes
\begin{itemize}\setlength\itemsep{3pt}
\item From country tax reports of International Bureau of Fiscal Documentation
\end{itemize}
\item Income tax base: \% EAP whose gross income is above income tax exemption threshold

\end{itemize}
\end{frame}

\begin{frame}
\frametitle{Fact1: Employee share $\Uparrow$ with income and development levels} 
\begin{columns}
\begin{column}{0.6\textwidth}
  \begin{figure}
    \includegraphics[keepaspectratio, width=\textwidth]{j2022_f1.png}
    \caption*{Red: Employee, Blue: Self-employed}
  \end{figure}
  
\end{column}  
\begin{column}{0.35\textwidth}
\begin{itemize}\setlength\itemsep{3pt}
\item Within country, employee share increase over income
\item Across countries, employee share profile moves leftward 
\item Over transition, increase in employee share concentrated in deciles gradually further down 
\end{itemize}  
\end{column}
\end{columns}
\end{frame}

\begin{frame}
\frametitle{Fact2: Exemption threshold moves downward as country develops} 
\begin{columns}
\begin{column}{0.6\textwidth}
  \begin{figure}
    \includegraphics[keepaspectratio, width=\textwidth]{j2022_f2.png}
    \caption*{Red: Employee, Blue: Self-employed, Vertical: Threshold}
  \end{figure}
  
\end{column}  
\begin{column}{0.35\textwidth}
\begin{itemize}\setlength\itemsep{3pt}
\item At low development, threshold at top deciles 
\item With development, growth in employee share comoves with falling thresholds
\item Tax base increases
\item Documents frequent updating of thresholds (base increase not driven by nominal inflation)
\end{itemize}  
\end{column}
\end{columns}
\end{frame}

\begin{frame}
\frametitle{Fact3: Threshold is located where employee share remains constant} 
\begin{columns}
\begin{column}{0.6\textwidth}
  \begin{figure}
    \includegraphics[keepaspectratio, width=\textwidth]{j2022_f3.png}
    \caption*{}
  \end{figure}
  
\end{column}  
\begin{column}{0.35\textwidth}
\begin{itemize}\setlength\itemsep{3pt}
\item Composition of tax base is constant and high in employee share
\item Employee share: 85-95\%
\item Suggest that employee composition matters in deciding threshold
\item \% Self-employed $\Uparrow\to$ Enforcement costs $\Uparrow$ (high employee share desirable)
\end{itemize}  
\end{column}
\end{columns}
\end{frame}

\begin{frame}
\frametitle{Takeaways from the stylized facts}
Cost channels at play in expanding tax base
\begin{itemize}\setlength\itemsep{3pt}
\item Concerns with taxing income segment with too high-self employed population
\begin{itemize}
\item Employees have information trails, self-employed do not
\end{itemize}
\item Compliance cost imposed on newly taxed self-employed may be high
\begin{itemize}
\item Employees have third-party withholding
\end{itemize}
\item Worsen between-group horizontal equity among employee vs self-employed
\begin{itemize}
\item Self-employed have evasion opportunities $\to$ effective taxes paid < legal liabilities 
\end{itemize}
\end{itemize}\medskip
Rise of modern taxation: Decrease in threshold (broad-based tax) happens when
\begin{itemize}\setlength\itemsep{3pt}
\item Increase in employee share at lower end of income distribution (Fact 1)
\item This lowers costs of reform and expand tax base (Fact 2)
\item While retaining high employee share, making tax administration manageable (Fact 3)
\end{itemize}
\end{frame}

\begin{frame}
\frametitle{Implications for developing countries} 
Redistribution by direct taxes might be ill-suited for developing countries
\begin{itemize}\setlength\itemsep{3pt}
\item Developed countries extensively use third-party information
\begin{itemize}
\item Evidenced by increase in employee share down the income distribution over time
\end{itemize}
\item Developing countries cannot enforce taxes flexibly due to cost constraints 
\begin{itemize}
\item Higher share of self-employment (difficult to gather info)
\end{itemize}
\item For such settings, indirect taxation could be the best option for redistribution
\end{itemize}\medskip
Remaining issues 
\begin{itemize}
\item How would implementation of technology affect tax base long-run?
\begin{itemize}\setlength\itemsep{3pt}
\item More use of third-party information $\to$ Rising `employee' share?
\item Expanding tax base would then be less costly $\to$ tax collection, public goods $\Uparrow$
\item But what about adoption (existence of political will)?
\end{itemize}
\end{itemize}
\end{frame}




\section{How tax administration affects tax collection: Evidence from experiments}
\subsection{Pomeranz (2015): No Taxation without Information: Deterrence and Self-Enforcement in the Value Added Tax}
\begin{transitionframe}
\begin{center}
  \begin{huge}
    \textcolor{NTHU1}{%
  Pomeranz (2015):\\ No Taxation without Information: Deterrence and Self-Enforcement in the Value Added Tax
    }
  \end{huge}
\end{center}
\end{transitionframe}



\begin{frame}
\frametitle{Q: How does third-party information in VAT help taxation?} 
VAT is believed to facilitate enforcement via a third-party information 
\begin{itemize}\setlength\itemsep{3pt}
\item Built-in incentives to ensure that transactions are revealed
\begin{itemize}\setlength\itemsep{3pt}
  \item Downstream firms demand receipt in order to deduct input costs from their liabilities
  \item This incentive is sustained until the transaction reaches the final consumer
\end{itemize}
\item Implementation: 47 countries in 1990 $\to$ 140 in 2010s 
\item But the effect of this built-in structure not empirically tested at this point
\end{itemize}\medskip
This paper uses randomized experiments to verify the properties of this structure
\begin{itemize}\setlength\itemsep{3pt}
  \item This experiment was conducted on more than 400,000 firms in Chile
  \item Letter message experiment to test preventive deterrence effect on evasion
  \item Spillover experiment to identify underlying mechanisms behind information effects
  \item Idea is to test how the paper trails interact with tax audit probabilities
\end{itemize}
\end{frame}



\begin{frame}
\frametitle{Conceptual framework of VAT} 
Setup
\begin{itemize}\setlength\itemsep{3pt}
\item Let $i\in\{1,2,...,N\}$ index the firm in a value chain where firm $i+1$ buys from firm $i$
\item $s_i$ and $c_i$ are the true sales and costs for firm $i$. $\widehat{s}_i$ and $\widehat{c}_i$ are reported sales and costs
\item Firm $i+1$ wants receipt from firm $i$ to deduct $s_i$ from its cost $c_{i+1}$ (No collusion here)
\item The last firm $N$ does not face this problem: Consumers do not require receipt
\end{itemize}\medskip
Types of potential evasion and role of perceived audit provbabilities
\begin{itemize}\setlength\itemsep{3pt}
  \item Unilateral: Firms underreport own liability while risking discrepancy with other firms
  \begin{itemize}\setlength\itemsep{3pt}
    \item $\widehat{c}_i>c_i$ or $\widehat{s}_i<s_i\implies\widehat{c}_i>\widehat{s}_{i-1}$ or $\widehat{s}_i < \widehat{c}_{i+1}$ 
  \end{itemize}
  \item Collusive: Both firms drop the sale entirely or underreport liabilities 
  \begin{itemize}\setlength\itemsep{3pt}
    \item $\widehat{c}_i \neq c_i$ but $\widehat{c}_i=\widehat{s}_{i-1}$ or $\widehat{s}_i \neq s_i$ but $\widehat{s}_i = \widehat{c}_{i+1}$ 
    \item If firm $n$ engage in this evasion, it must be the case for all $i\geq n$ up to $i=N$ 
    \item It must also be that $\widehat{s}_i-\widehat{c}_i \leq s_i-c_i \iff c_i-\widehat{c}_i \leq s_i - \widehat{s}_i$
  \end{itemize}
  \item As audits can uncover paper trail, the threat of audit should deter evasion
\end{itemize}
\end{frame}

\begin{frame}
  \frametitle{Prediction on tax evasion behaviors after announcement of audit}
  Unilateral evasion
\begin{itemize}\setlength\itemsep{3pt}
\item Suppose audit is announced on firm $n$ in the middle of the chain
\item This forces supplier and client firms to revise perceived audit probability
\item Thus, they revise declared transactions to be consistent with firm $n$
\begin{itemize}\setlength\itemsep{3pt}
  \item Supplier: Increase sales figures and VAT payment
  \item Client: Decrease cost figures and increase VAT payment
\end{itemize}
\end{itemize}\medskip
Collusive evasion
\begin{itemize}\setlength\itemsep{3pt}
\item Again, firm $n$ receives audit announcement
\item Firm $n$ would be incentivized to create a paper trail (increased sale and costs)
\item Supplier firm induced to increase reported sales and VAT payment 
\begin{itemize}
\item Possible chain reaction continues to onto its own supplier
\end{itemize}
\item Client firms has room to increase their cost and reduce VAT 
\begin{itemize}
\item Since supplier raised sale, client has more room to declare its costs
\end{itemize}
\end{itemize}
\end{frame}

\begin{frame}
\frametitle{Experiment design} 
Experiment 1: Letter experiment to check behavior of a single firm
\begin{itemize}\setlength\itemsep{3pt}
\item Deterrence: Reminds firms of the possilbility of audit (102,031)
\item Tax Morale: Highlight social norms on taxes, but no content on deterrence (18,579)
\item Placebo: New feature on tax website to see if messaging alone has impacts (18,519)
\item Outcome variable: VAT filings of each firms
\end{itemize}\medskip
Experiment 2: Spillover experiment to see if tax compliance on other firms are affected
\begin{itemize}\setlength\itemsep{3pt}
  \item Sample of 5,600 suspected firms already chosen for audit
  \item Randomly pre-announce audit to half of the firms
  \item Data on partner firms collected in the audit
  \item Check VAT declarations before and after announcemenrt
  \end{itemize}
\end{frame}

\begin{frame}
  \frametitle{Deterrence message increases compliance (reduce evasion)}
  \begin{figure}
    \centering
    \includegraphics[keepaspectratio, width=0.8\textwidth]{p2015_t1.png}
  \end{figure}
\end{frame}

\begin{frame}
  \frametitle{Preannouncement also affects partner firms' taxation response}
  \begin{figure}
    \centering
    \includegraphics[keepaspectratio, width=0.7\textwidth]{p2015_t2.png}
  \end{figure}
\end{frame}

\begin{frame}
\frametitle{Key takeaways and discussion} 
Key findings
\begin{itemize}\setlength\itemsep{3pt}
\item Paper trail in VAT self-enforces accurate tax information
\item Transactions already under paper trail more likely to comply with audit possibility
\item Increasing audit probability increases tax payment of supplier firms
\begin{itemize}
  \item Not necessarily downstream firms: Consumers have less incentives to ask for receipt 
\end{itemize} 
\end{itemize}\medskip
Overall message
\begin{itemize}\setlength\itemsep{3pt}
  \item Interaction of information and enforcement matters
  \item Information from paper trail can be a powerful tool
   \begin{itemize}
    \item Either in form of VAT or digitalized property rasters or earnings 
   \end{itemize}
  \item But it also explains why tax collection in developing countries are low
   \begin{itemize}\setlength\itemsep{3pt}
    \item Many home-based activities with no paper trail
    \item Lack of credible threat/perception of enforcement
   \end{itemize}
  \end{itemize}
\end{frame}

\subsection{De Neve, Imbert, Spinnewijn, Tsankova, Lutz (2021): How to Improve Tax Compliance? Evidence from Population-Wide Experiments in Belgium}
\begin{transitionframe}
\begin{center}
  \begin{huge}
    \textcolor{NTHU1}{%
De Neve, Imbert, Spinnewijn, Tsankova, Lutz (2021): \\ How to Improve Tax Compliance? Evidence from Population-Wide Experiments in Belgium
    }
  \end{huge}
\end{center}
\end{transitionframe}


\begin{frame}
\frametitle{Q: How do  different administrative approaches affect tax compliance?} 
Tax compliance matters for healthy functioning societies
\begin{itemize}\setlength\itemsep{3pt}
 \item Cost of noncompliance in US: 319 billion USD annually
  \item Truthfulness and timely payment both matters for compliance
\item Increasing yet incomplete understanding on nonpecuniary motives for tax compliance 
\begin{itemize}
  \item Recent interest in `tax morale'
\end{itemize} 
\end{itemize}\medskip
This paper studies how communications from tax authorities affects tax compliance
\begin{itemize}\setlength\itemsep{3pt}
  \item Different dimensions: Simplified communication vs deterrence vs appeal to tax morale
  \item Messaging experiment from tax authorities in Belgium (2014-16)
  \item Estimates how different dimensions of tax reminders differentially affect compliance
  \begin{itemize}\setlength\itemsep{3pt}
  \item Message content
  \item All taxpayers vs taxpayers who missed deadline
\end{itemize} 
  \end{itemize}
\end{frame}



\begin{frame}
\frametitle{Tax process in Belgium} 
\begin{figure}
  \centering
  \includegraphics[keepaspectratio, width=0.6\textwidth]{distl2021_f1.png}
\end{figure}

\end{frame}

\begin{frame}
\frametitle{Experiment design} 
Treatment
\begin{itemize}\setlength\itemsep{3pt}
 \item Simplification: Shorter letters, reduce information overload, action-relevnt info
 \item Deterrence: Simplified + make penalities explicit
 \item Tax deterrence: Simplified + emphasize social benefits of taxation
 \item Control group: Standard long message with complicated information
\end{itemize}\medskip
Implementation of the experiment
\begin{itemize}\setlength\itemsep{3pt}
  \item Four experiments: Filing, Filing reminder, Payment, Payment reminder
  \begin{itemize}\setlength\itemsep{3pt}
    \item Filing experiment only had tax morale
    \item Filing reminder experiment did not have morale component
  \end{itemize}
  \item Only the ones who missed the deadlines are part of the reminder experiment
  \item Randomized based on birthday (payment) and national ID number (other experiments)
  \item Regression is an OLS where
  \[
  y_i =\alpha + \beta_s S_i + \sum_j \beta_j T^j_i+\gamma X_i + e_i
  \]
  \end{itemize}
\end{frame}

\begin{frame}
\frametitle{Different effects on tax payment} 
\begin{figure}
  \centering
  \includegraphics[keepaspectratio, width=0.6\textwidth]{distl2021_f2.png}
\end{figure}
\end{frame}

\begin{frame}
\frametitle{Exploring mechanisms} 
Simplication works primarily with extensive margins
\begin{itemize}\setlength\itemsep{3pt}
  \item Make relevant deadline salient: Effects larger around the deadlines
  \item Does not lead to large increases in the amount of taxes paid
\end{itemize}\medskip
Deterrence component complements the effect of simplified messages
\begin{itemize}\setlength\itemsep{3pt}
  \item Makes deadline and relevant information salient and encourage immediate action
  \item No large effect on the amount of taxes paid
  \item Tax morale part may have increased relevant knowledge but ineffective beyond this
\end{itemize}\medskip
Heterogeneity
\begin{itemize}\setlength\itemsep{3pt}
  \item Simplication more effective on households with children and median solvency score
  \item Deterrence most effective on young taxpayers, and those with low liabilities
\end{itemize}
\end{frame}

\begin{frame}
\frametitle{Cost-effectiveness of interventions and remaining discussion} 
Cost effectiveness of different interventiones 
\begin{itemize}\setlength\itemsep{3pt}
  \item The authors estimate the costs using RD estimates around enforcement intensity
  \item Based on this cost-effectiveness is measured in three ways 
  \item In all of these results: Costs of simplification $<$ Costs of extra enforcement
  \begin{itemize}\setlength\itemsep{3pt}
    \item Simplification: Total cost is 10\% of extra deterrence measure in one of the approaches
  \end{itemize} 
\end{itemize}\medskip
Discussions
\begin{itemize}\setlength\itemsep{3pt}
  \item Compliance can increase if directions are straightforward
  \item Morale message not effective: Likely because people are aware of importance of taxes
  \item Lack of compliance, thus, is largely explained by inattention of taxpayers 
  \item Can this apply on tax collectors?
  \begin{itemize}\setlength\itemsep{3pt}
  \item They know that as public servants, collecting taxes are crucial
  \item But they may be techinically unaware of deviant behaviors...
\end{itemize}
\end{itemize}
\end{frame}


\section{Effect of tax collectors on compliance}
\subsection{Basri, Felix, Hanna, Olken (2021): Tax Administration versus Tax Rates: Evidence
from Corporate Taxation in Indonesia}
\begin{transitionframe}
\begin{center}
  \begin{huge}
    \textcolor{NTHU1}{%
Basri, Felix, Hanna, Olken (2021):\\ Tax Administration versus Tax Rates: Evidence
from Corporate Taxation in Indonesia
    }
  \end{huge}
\end{center}
\end{transitionframe}

\begin{frame}
\frametitle{Q: How much does improving tax adminstration raise tax revenues?}
Tax revenues in developing countries remain low
\begin{itemize}\setlength\itemsep{3pt}
\item Is it due to high elasticity of taxable income (ETI)?
\begin{itemize}\setlength\itemsep{3pt}
\item High ETI implies that tax revenues fall in large amounts w.r.t. rise in tax rates
\end{itemize}
\item Or is it a problem of tax administration?
\begin{itemize}\setlength\itemsep{3pt}
\item Enhanced administration may make evasion/avoidance much more difficult
\item Rise in tax revenues without rise in tax rates
\end{itemize}
\end{itemize}\medskip
This paper studies this question in context of corporate taxes in Indonesia
\begin{itemize}\setlength\itemsep{3pt}
\item Administrative reforms created medisum size taxpayer offices (MTO)
\begin{itemize}\setlength\itemsep{3pt}
\item higher ($\times3$) staff-to-taxpayer ratio than other offices to increase enforcement
\end{itemize}
\item First study the effect of intensified administration in administrative tax data
\item Then compare with changes in statutory corporate income tax rates several yrs later
\end{itemize}
\end{frame}


\begin{frame}
\frametitle{Relevant tax reforms in Indonesia}
Administrative reforms in Indonesia in mid-2000s
\begin{itemize}\setlength\itemsep{3pt}
\item Moved several hundred top corporate taxpayers across 19 tax regions to MTOs 
\item MTOs exclusively focused on large taxpayers and had more staff than other offices
\item Taxpayers not assigned to MTOs assigned to smaller offices (PTOs)
\item Staffs in MTOs and PTOs are similar in experience and performance assessments
\end{itemize}\medskip
Tax rate reform on corporate income tax
\begin{itemize}\setlength\itemsep{3pt}
\item In 2009, Indonesia changed from progressive 15-30\% rate to flat 28\% tax rate
\item Pre-reform tax base was net taxable profits
\item Post-reform taxes are determined based on gross revenues
\item Total tax and discount determined by a set formula across 3 revenue brackets
\end{itemize}
\end{frame}


\begin{frame}
\frametitle{Data sources and regression}
Data sources
\begin{itemize}\setlength\itemsep{3pt}
\item Anonymized microdata on corporate taxpayers (2003-2011)
\item Detail on tax reporting, wage bills, audits, VAT assessments, and tax payment
\end{itemize}\medskip
Regression approach
\begin{itemize}\setlength\itemsep{3pt}
\item MTO reform: Use matched DID
\[
y_{it} = \alpha + \beta (M_{i}\times 1[t>2005])+ \delta_t + \delta_i + e_{it}
\]\vspace{-20pt}
\begin{itemize}\setlength\itemsep{3pt}
  \item Control: Taxpayers in same region with similar pre-MTO tax payments and gross revenue
  \item Match control and treatment based on entropy balancing (Hainmueller 2012)
\end{itemize}
\item Effects of tax rates: Instrument change in tax rates with tax formula following reform
\begin{itemize}\setlength\itemsep{3pt}
  \item The goal is to estimate elasticity of taxable income w.r.t tax rate changes
  \item Lower elasticity implies that taxable income does not respond much to tax rate
  \item Used to calculate how much tax rate should change to match revenue from MTOs
\end{itemize}
\end{itemize}
\end{frame}


\begin{frame}
  \frametitle{Tax revenues following MTO nearly doubles}
  \begin{figure}
    \centering
    \includegraphics[keepaspectratio, width=1\textwidth]{bfho2021_f1.png}
  \end{figure}
\end{frame}

\begin{frame}
  \frametitle{Effects driven by higher reported incomes}
  \begin{figure}
    \centering
    \includegraphics[keepaspectratio, width=1\textwidth]{bfho2021_f2.png}
  \end{figure}
\end{frame}

\begin{frame}
  \frametitle{ETIs are actually similar across different firms}
  \begin{figure}
    \centering
    \includegraphics[keepaspectratio, width=0.6\textwidth]{bfho2021_t1.png}
  \end{figure}
  \begin{itemize}\setlength\itemsep{3pt}
    \item ETI is 0.579, similar to many developed countries
    \item ETIs statistically similar = MTOs did not work by reducing ETIs
  \end{itemize}
\end{frame}

\begin{frame}
  \frametitle{Other key findings and discussion}
  Other key findings
  \begin{itemize}\setlength\itemsep{3pt}
    \item Tax enforcement no longer size-dependent
    \begin{itemize}
      \item Before: Small offices focusing on large firms 
      \item This created disincentives for small firms to grow
      \item Now: MTOs able to treat firms of various sizes equally 
    \end{itemize}
    \item Tax revenue increases by administrative reform can be cost-effective
     \begin{itemize}\setlength\itemsep{3pt}
     \item Improved tax administration = raising top rates on all firms by 8 percentage points
     \item So improving administration can have significant returns for developing countries
    \end{itemize}
  \end{itemize}\medskip
  Discussion: How to improve tax administration
   \begin{itemize}\setlength\itemsep{3pt}
     \item Improving staff-to-taxpayer ratio by raising staffing
     \item Alternatively, specialized offices may help
     \item Another dimension is the quality of staffing (education, capital, training)
     \begin{itemize}\setlength\itemsep{3pt}
     \item Better trained, well-equipped staffing can improve effective number of personnel
     \end{itemize}
    \end{itemize}
\end{frame}


\subsection{Balán, Bergeron, Tourek, Weigel (2022): Local Elites as State Capacity: How City Chiefs Use Local Information to Increase Tax Compliance in the Democratic Republic of Congo}
\begin{transitionframe}
\begin{center}
  \begin{huge}
    \textcolor{NTHU1}{%
Balán, Bergeron, Tourek, Weigel (2022):\\ Local Elites as State Capacity: How City Chiefs Use Local Information to Increase Tax Compliance in the Democratic Republic of Congo
    }
  \end{huge}
\end{center}
\end{transitionframe}

 
\begin{frame}
\frametitle{Q: How can local information be used effectively to increase taxes?} 
How can fiscal capacities develop in presence of local \& traditional elites
\begin{itemize}\setlength\itemsep{3pt}
\item On the one hand, can capture politics and civic society
\item On the other, could be room to collaborate to improve governance
\item Potential trade-offs
\begin{itemize}
\item Enforcement backed by the state vs local knowledge and coordination?
\end{itemize}
\end{itemize}\medskip
This study investigates if local elites can contribute to statebuilding
\begin{itemize}\setlength\itemsep{3pt}
\item Comparison between state agents vs local elites
\begin{itemize}\setlength\itemsep{3pt}
\item Local elites know more contextual knowledge but could lead to leakage in revenues
\item State agents have better enforcement capacity
\end{itemize}
\item Based on field experiment randomizing type of tax collectors across neighborhoods
\item Estimate changes in revenues and bribes, with tests for various mechanism
\end{itemize}
\end{frame}

\begin{frame}
\frametitle{Taxation in Democratic Republic of Congo (DRC)} 
Setting of DRC
\begin{itemize}\setlength\itemsep{3pt}
\item Low-capacity state: tax/GDP ranks 188 out of 200 and  \$0.30 per capita tax revenue
\item Relies on property tax: First campaign to expand property tax in 2016
\item Public goods and services are scarce and low-quality
\begin{itemize}
\item Limited running water, electricity, unrepaired roads etc
\end{itemize}
\end{itemize}\medskip
The current study focuses on the second campaign (May-Dec 2018)
\begin{itemize}\setlength\itemsep{3pt}
\item Collectors trained by provincial ministry, separate for state and chiefs
\item Involved property registration updates, providing liability information 
\item Followed by tax visits (revisit until taxes are paid)
\item Collectors compensated based on houses registered and revenue collection
\item Property tax rates are tiered flat, fixed fees (low: \$2 tax, 89\% vs high: \$9 tax, 11\%)
\begin{itemize}
\item Equivalent to 0.32\% of property value ($\simeq$ lower end of property tax rates in US)
\end{itemize}
\end{itemize}\medskip
\end{frame}

\begin{frame}
\frametitle{Data and empirical strategy} 
RCT on 45000+ properties in 356 neighborhoods in Kasai-Central
\begin{itemize}\setlength\itemsep{3pt}
\item Randomized at neighborhood, treatment at past campaigns, experience of chiefs
\item Central (state agents) vs Local (chiefs) vs State agents w/ local info vs state and local
\item Pure control group: Declare their property values voluntarily
\item Other aspects of collection equal
\begin{itemize}
\item Property registration, assessment, tax liabilities, training, collection compensation
\end{itemize}
\item Complemented with tax administrative data and household surveys
\end{itemize}\medskip
Estimation leverages random assignment of treatment (indiv $i$, neighbor $j$, strata $k$, time $t$)
\[
y_{ijkt}=\beta_0 + \beta\text{Local}_{jkt}+\gamma X_{ijkt} + \alpha_k + \theta_t + e_{ijkt}
\]\vspace{-20pt}
\begin{itemize}\setlength\itemsep{3pt}
\item For mechanisms, other treatment arms included as separate variable
\end{itemize}
\end{frame}

\begin{frame}
\frametitle{Increase in compliance and revenues} 
\begin{figure}
  \includegraphics[keepaspectratio, width=0.8\textwidth]{bbtw2022_t1.png}
\end{figure}
\end{frame}

\begin{frame}
\frametitle{Accurate assessments \& trust in government $\uparrow$, but bribes $\uparrow$} 
\begin{figure}
  \includegraphics[keepaspectratio, width=0.7\textwidth]{bbtw2022_t2.png}
\end{figure}
\end{frame}

\begin{frame}
\frametitle{Evidence of local information at play} 
\begin{figure}
  \includegraphics[keepaspectratio, width=1\textwidth]{bbtw2022_t3.png}
\end{figure}
\end{frame}

\begin{frame}
\frametitle{Discussion and takeaways}
What explains the chiefs' advantage? Better targeting with information advantage
\begin{itemize}\setlength\itemsep{3pt}
\item Lower effort cost? Chiefs did not make more visits
\item Target their visits? When state agents received info from chiefs, compliance $\Uparrow$
\item Better persuasion? Neutralize info advantage (visit in linear route)$\Rightarrow$ Central = Local
\end{itemize}\medskip
Takeaways and remaining questions
\begin{itemize}\setlength\itemsep{3pt}
\item Should countries with low capacities rely on local elites?
\begin{itemize}\setlength\itemsep{3pt}
\item Raise revenues, but also bribes
\item To shift to state, value of \$1 lost in bribes $>$$>$ \$1 gained in revenues
\end{itemize}
\item For technology to increase taxes, need to surpass some state capacity threshold?
\begin{itemize}\setlength\itemsep{3pt}
\item Speak to the complementary/substitutional relationship between tech and person
\end{itemize}
\end{itemize}\medskip
\end{frame}

 
\begin{frame}
\frametitle{Unanswered questions} 
Effects of technological advancements on taxation
\begin{itemize}\setlength\itemsep{3pt}
\item More countries adopting digital registries, automated paper trails, and identification
\item Many evidence that it led to better detection and compliance
\item Other questions that remain not yet fully explored 
 \begin{itemize}\setlength\itemsep{3pt}
\item Can it lead to better training of tax officials?
\item Can it lead to change in perception about taxation among citizens?
\end{itemize}
\end{itemize}\medskip
What can better administration of taxes do
\begin{itemize}\setlength\itemsep{3pt}
\item More efforts to increase administration in place
\item But what about administration for how to \textit{use} these taxes?
 \begin{itemize}\setlength\itemsep{3pt}
\item How can developing countries provide public services efficiently with these taxes?
\item How to ensure that these officials appropriately distribute taxes without corruption?
\end{itemize}
\item We will cover topics on public service and personnel in developing countries 
\end{itemize}
\end{frame}
%%%%%%%%%%%
\end{document}
